The observation table shows that the measured TLD counts decrease as the radial distance from the source increases. This is the expected physical behavior of brachytherapy photons, since the photon fluence is reduced by geometric spreading, attenuation in the phantom medium, and the finite source size. The highest response is obtained near the reference transverse plane, $\theta = 90^\circ$, where the path through the source encapsulation is minimal and the geometry is most favorable.

The variation of counts with angle also indicates the anisotropic nature of the source. Angles on opposite sides of the source show nearly symmetric behavior, while the response away from the transverse direction changes because of self-absorption inside the active core and filtering by the capsule. The sign change in $\beta$ only reflects the orientation of the line joining the source end points to the point of interest; physically, the relevant quantity is the magnitude of the angular span. After geometry correction, the reference direction remains normalized to unity, confirming the TG-43 behavior of the source and validating the use of the line-source formalism for in-phantom dosimetry.

More generally, these dosimetric parameters are controlled by a combination of physical and geometric factors. The absorbed dose rate decreases with distance because of the inverse square law and additional attenuation in the phantom medium. The finite active length of the source introduces line-source geometry, so the dose depends not only on $r$ but also on the polar angle $\theta$ through the angular span $\beta$. The encapsulation material, source self-absorption, and photon scatter in the surrounding medium all contribute to the radial dose function and anisotropy function. In the transverse plane, the geometry is most symmetric, so the measured response is used as the reference condition for normalizing the TG-43 parameters. These factors together explain the observed variation in TLD counts and the physical significance of the measured dosimetry trends.

For the anisotropy function, the theoretical expectation is that $F(r,\theta)$ is maximum at $\theta = 90^\circ$ and decreases as the angle moves away from the transverse plane. In the present experimental data, this decreasing pattern is not perfectly visible at all points. This small irregularity is likely due to random and experimental uncertainties, such as TLD positioning uncertainty, readout fluctuation, finite detector size, and minor setup misalignment during irradiation.

\begin{figure}[H]
\centering
\includegraphics[width=0.9\textwidth]{\detokenize{Figures/Theory/MonteCarlo.png}}
\caption{Comparison of the anisotropy function, F(r,u ), values published by several authors, estimated either experimentally or using MC simulations for the microSelectron HDR 192Ir type old ``classic'' source at radial distance (a) r= 3.0 cm and (b) r = 5.0 cm.}
\end{figure}